\documentclass[11pt,a4paper]{article}
\usepackage[margin=2.5cm]{geometry}
\usepackage{amsmath,amssymb}
\usepackage{booktabs}
\usepackage{listings}
\usepackage{xcolor}
\usepackage{hyperref}
\usepackage{graphicx}
\usepackage{array}
\usepackage{multirow}
\usepackage{caption}
\usepackage{float}
\usepackage[utf8]{inputenc}
\usepackage[T5]{fontenc}
\usepackage[vietnamese]{babel}
\usepackage{enumitem}
\usepackage{tcolorbox}
\tcbuselibrary{skins,breakable}

% Code listing style
\lstset{
  language=Matlab,
  basicstyle=\ttfamily\footnotesize,
  keywordstyle=\color{blue},
  commentstyle=\color{gray!70},
  stringstyle=\color{orange!80!black},
  frame=single,
  breaklines=true,
  numbers=left,
  numberstyle=\tiny\color{gray},
  xleftmargin=1.5em,
  framexleftmargin=1.5em,
}

\hypersetup{colorlinks=true, linkcolor=blue, citecolor=blue, urlcolor=blue}

\title{
  \Large\textbf{Cài đặt Mode B --- Refined Type I Single-Panel Codebook}\\[0.4em]
  \large TS~38.214 \S5.2.2.2.1a (3GPP Release~19)\\[0.4em]
  \normalsize Lý thuyết, Thuật toán và Tích hợp vào MATLAB 5G Toolbox
}
\author{Thang T.Q.}
\date{Tháng 3, 2026}

\begin{document}
\maketitle
\tableofcontents
\newpage

%% ============================================================
\section{Tóm tắt}
%% ============================================================

Báo cáo này ghi chép chi tiết toàn bộ quá trình cài đặt \textbf{Mode B} của
Refined Codebook Type~I Single-Panel (TS~38.214 \S5.2.2.2.1a, Rel-19)
vào thư viện MATLAB tùy chỉnh.
Đây là phần tiếp nối sau khi Mode~A đã được hoàn thiện trong báo cáo trước
(\texttt{report\_rel19\_upgrade\_vi.pdf}).

\medskip
\noindent\textbf{Những gì đã hoàn thành:}
\begin{itemize}
  \item Phân tích lý thuyết cấu trúc codebook Mode~B theo tiêu chuẩn.
  \item Nhận ra vấn đề kích thước không gian tìm kiếm (4,3~GB ở hạng~2) và
        thiết kế thuật toán \emph{Greedy Factored Search} để thay thế.
  \item Cài đặt hàm \texttt{getPMIType1SinglePanelR19CodebookModeB()}
        (246~dòng) và tích hợp vào \texttt{nrPMIReport.m} qua 5 điểm sửa đổi.
  \item Viết script kiểm tra \texttt{test\_ModeB\_verification.m} với 5 bài test
        --- tất cả đều PASS, bao gồm chứng minh
        $\text{cap}(\text{Mode~B}) \geq \text{cap}(\text{Mode~A})$.
\end{itemize}

%% ============================================================
\section{Lý thuyết: Mode B theo TS~38.214 \S5.2.2.2.1a}
%% ============================================================

\subsection{Khác biệt cơ bản giữa Mode A và Mode B}

Cả hai mode đều xây dựng bộ tiền mã hoá dạng dual-polarisation DFT:
\begin{equation}
  W^{(v)} = \frac{1}{\sqrt{v \cdot P}}
  \begin{bmatrix}
    v_{l_1,m_1} & \cdots & v_{l_v,m_v} \\
    \varphi(c_1)\,v_{l_1,m_1} & \cdots & \varphi(c_v)\,v_{l_v,m_v}
  \end{bmatrix}
  \label{eq:precoder}
\end{equation}
với $P = 128$, $\varphi(c) = e^{j\pi c/2}$, $c \in \{0,1,2,3\}$.

\medskip
Sự khác biệt nằm ở cách \textbf{xác định chỉ số beam} $(l_k, m_k)$:

\begin{table}[H]
\centering
\caption{So sánh cấu trúc chỉ số Mode A vs Mode B (hạng~$v$)}
\begin{tabular}{p{3cm}p{5.5cm}p{5.5cm}}
\toprule
 & \textbf{Mode A} & \textbf{Mode B} \\
\midrule
Chỉ số $i_{1,1}$ & $l \in \{0,\ldots,N_1 O_1{-}1\}$ (beam hàng) & $(q_1, q_2)$: offset phân số chung \\
Chỉ số $i_{1,2}$ & $m \in \{0,\ldots,N_2 O_2{-}1\}$ (beam cột) & $i_{1,2,l} \in \{0,\ldots,N_1 N_2{-}1\}$ per-layer \\
Beam thứ $k>1$ & offset cố định $(k_1,k_2)$ từ beam 1 & độc lập, chọn tự do \\
Tổng số ứng viên (hạng 1) & $N_1 O_1 \times N_2 O_2 \times 4 = 4096$ & $O_1 O_2 \times N_1 N_2 \times 4 = 4096$ \\
Tổng số ứng viên (hạng 2) & $4096 \times 4 = 16{,}384$ & $16 \times 64^2 \times 16 \approx 1\text{M}$ \\
\bottomrule
\end{tabular}
\end{table}

\noindent\textbf{Ưu điểm của Mode~B}: mỗi tầng được phép chọn một DFT beam
hoàn toàn độc lập. Đây là lợi thế lớn trên các kênh nhiều tầng
(multi-rank) có singular vectors phân tán, ví dụ CDL-B.

\subsection{Ánh xạ chỉ số Mode B (Bảng 5.2.2.2.1a-6/7)}

Với cấu hình $(N_1, N_2) = (16, 4)$, $(O_1, O_2) = (4, 4)$:

\begin{enumerate}
  \item \textbf{Offset phân số}: $i_{1,1} = [q_1, q_2]$, với
        $q_1 \in \{0,\ldots,O_1{-}1\}$ và $q_2 \in \{0,\ldots,O_2{-}1\}$.
        Mỗi cặp $(q_1, q_2)$ xác định một \emph{lưới con DFT} trong tổng lưới
        $N_1 O_1 \times N_2 O_2$.
  \item \textbf{Chỉ số beam theo tầng}: cho tầng $l$, $i_{1,2,l} \in \{0,\ldots,N_1 N_2{-}1\}$.
        Ánh xạ từ chỉ số $i_{1,2,l}$ sang tọa độ lưới con:
        \begin{align}
          n &= N_1 N_2 - 1 - i_{1,2,l} \label{eq:n_mapping}\\
          n_1 &= n \bmod N_1, \quad n_2 = \lfloor n / N_1 \rfloor \label{eq:n1n2}\\
          m_1 &= O_1 n_1 + q_1, \quad m_2 = O_2 n_2 + q_2 \label{eq:m1m2}
        \end{align}
  \item \textbf{Vector beam DFT}:
        \begin{equation}
          v_{m_1,m_2} = \mathbf{u}_h(m_1) \otimes \mathbf{u}_v(m_2), \quad
          [\mathbf{u}_h(m_1)]_k = e^{j2\pi m_1 k/(O_1 N_1)}, \quad
          [\mathbf{u}_v(m_2)]_k = e^{j2\pi m_2 k/(O_2 N_2)}
          \label{eq:vlm}
        \end{equation}
  \item \textbf{Co-phase}: $i_2 = [c_1,\ldots,c_v]$, $c_l \in \{0,1,2,3\}$,
        áp dụng $\varphi(c_l)$ cho phân cực dưới của tầng $l$.
\end{enumerate}

\subsection{Cấu trúc bộ tiền mã hoá Mode B (Bảng 5.2.2.2.1a-7)}

Cho hạng $v$ ($v = 1,2,3,4$), bộ tiền mã hoá là:
\begin{equation}
  W^{(v)} = \frac{1}{\sqrt{v \cdot P}}
  \underbrace{\begin{bmatrix}
    \mathbf{V}_{\text{top}} \\
    \boldsymbol{\Phi}\,\mathbf{V}_{\text{top}}
  \end{bmatrix}}_{P \times v}, \quad
  \mathbf{V}_{\text{top}} =
  \bigl[v_{m_1^{(1)},m_2^{(1)}},\,\ldots,\,v_{m_1^{(v)},m_2^{(v)}}\bigr],
  \quad
  \boldsymbol{\Phi} = \mathrm{diag}\bigl(\varphi(c_1),\ldots,\varphi(c_v)\bigr)
  \label{eq:W_modeB}
\end{equation}

\noindent\textbf{Kiểm tra chuẩn hoá}:
\begin{equation}
  \|W^{(v)}\|_F^2 = \frac{1}{v P} \cdot
  \bigl(\|\mathbf{V}_{\text{top}}\|_F^2 + \|\boldsymbol{\Phi}\,\mathbf{V}_{\text{top}}\|_F^2\bigr)
  = \frac{1}{vP} \cdot 2 \sum_{l=1}^{v} \|v_{m_1^{(l)},m_2^{(l)}}\|^2
  = \frac{1}{vP} \cdot 2vP/2 = 1
  \label{eq:norm_check}
\end{equation}
vì mỗi $\|v_{l,m}\|^2 = N_1 N_2 = P/2$.

\subsection{Lưới DFT Mode B bao phủ cùng không gian với Mode A}

Nhận xét quan trọng: với hạng~1, không gian tìm kiếm Mode~B
\emph{bao trùm chính xác} không gian của Mode~A.
Mỗi cặp $(q_1, q_2) \in \{0,\ldots,O_1{-}1\}\times\{0,\ldots,O_2{-}1\}$
tạo ra một lưới con $\{m_1 = O_1 n_1 + q_1\} \times \{m_2 = O_2 n_2 + q_2\}$,
và hợp của $O_1 O_2 = 16$ lưới con đó chính là toàn bộ lưới
$\{0,\ldots,N_1 O_1{-}1\} \times \{0,\ldots,N_2 O_2{-}1\}$ mà Mode~A sử dụng.
Điều này đảm bảo:
\begin{equation}
  \text{cap}(W^{(1)}_\text{ModeB}) = \text{cap}(W^{(1)}_\text{ModeA})
  \quad \text{(hạng 1, khi tìm kiếm đầy đủ)}
  \label{eq:equiv_rank1}
\end{equation}

%% ============================================================
\section{Thách thức Kỹ thuật và Phân tích Thiết kế}
%% ============================================================

\subsection{Vấn đề Kích thước Codebook}

Với Mode~A, codebook được \emph{tính trước} (precomputed) dưới dạng tensor
$\mathbf{C} \in \mathbb{C}^{P \times v \times N_{i_2} \times N_{i_{1,1}} \times N_{i_{1,2}} \times N_{i_{1,3}}}$
rồi truyền vào hàm tìm PMI. Cách tiếp cận này không áp dụng được cho Mode~B
vì số ứng viên codebook tăng theo lũy thừa với hạng:

\begin{table}[H]
\centering
\caption{Kích thước không gian tìm kiếm Mode B}
\begin{tabular}{ccccc}
\toprule
Hạng $v$ & $O_1 O_2$ & $N_1 N_2$ (beam/tầng) & $4^v$ (co-phase) & Tổng ứng viên \\
\midrule
1 & 16 & 64 & 4 & 4\,096 \\
2 & 16 & $64^2 = 4{,}096$ & 16 & 1\,048\,576 \\
3 & 16 & $64^3 = 262{,}144$ & 64 & $\approx 2{,}7 \times 10^8$ \\
4 & 16 & $64^4 = 16{,}777{,}216$ & 256 & $\approx 6{,}9 \times 10^{10}$ \\
\midrule
\multicolumn{4}{l}{Bộ nhớ tensor hạng 2 (complex128)} & $\approx 4{,}3$~GB \\
\bottomrule
\end{tabular}
\end{table}

\noindent Hạng~2 đã yêu cầu 4,3~GB RAM --- \textbf{không khả thi} để tính trước.
Do đó cần thiết kế thuật toán tìm kiếm trực tuyến.

\subsection{Chiến lược: Greedy Factored Search}

Thay vì tìm kiếm vét cạn toàn bộ không gian, thuật toán được chia thành
ba giai đoạn phân cấp:

\begin{tcolorbox}[colback=blue!5,colframe=blue!50,title=Greedy Factored Search --- Ba giai đoạn]
\begin{description}[leftmargin=2em]
  \item[Giai đoạn 1:] Liệt kê $O_1 O_2 = 16$ cặp offset $(q_1, q_2)$.
  \item[Giai đoạn 2:] Với mỗi $(q_1, q_2)$:
    \begin{itemize}
      \item \textbf{Hạng 1}: duyệt tất cả $N_1 N_2 = 64$ beam (vét cạn hoàn toàn).
      \item \textbf{Hạng $v > 1$}: dùng SVD wideband $\mathbf{H}_\text{wb}$ để chọn 1 beam tốt nhất
            cho mỗi tầng (greedy, phức tạp tuyến tính).
    \end{itemize}
  \item[Giai đoạn 3:] Duyệt vét cạn $4^v$ tổ hợp co-phase cho bộ beam đã chọn;
        chọn $W$ tối đa hoá dung lượng Shannon $\log_2\det(\mathbf{I} + \frac{1}{\sigma^2}\mathbf{H}_\text{wb}\mathbf{W}\mathbf{W}^H\mathbf{H}_\text{wb}^H)$.
\end{description}
\end{tcolorbox}

\medskip
\noindent\textbf{Độ phức tạp so sánh}:
\begin{itemize}
  \item Hạng 1: $16 \times 64 \times 4 = 4{,}096$ đánh giá (bằng Mode~A, đảm bảo tối ưu).
  \item Hạng 2: $16 \times (64 + 16) = 1{,}280$ đánh giá tổng (so với $> 10^6$ vét cạn).
  \item Hạng 4: $16 \times (64 + 256) = 5{,}120$ đánh giá tổng.
\end{itemize}

\subsection{Vấn đề Chọn Beam bằng SVD: Cả hai phân cực}

Trong cài đặt ban đầu, việc chấm điểm beam cho hạng $v > 1$ dùng hình chiếu lên
\emph{nửa trên} của vector kỳ dị phải:
\begin{equation}
  \text{score}(n,l)_\text{sai} = |V_q^H \mathbf{V}_\text{top}|_{n,l}^2
  \quad \text{(chỉ phân cực trên)}
\end{equation}
Điều này bỏ qua phân cực dưới. Với kênh Rayleigh ngẫu nhiên, $\mathbf{H}_\text{top}$
và $\mathbf{H}_\text{bot}$ độc lập nên vector kỳ dị $\mathbf{v}_l \in \mathbb{C}^{2N_1N_2}$
phân bổ năng lượng đồng đều giữa hai nửa. Hình chiếu chỉ lên nửa trên
bỏ mất $\approx 50\%$ thông tin.

\medskip\noindent\textbf{Sửa chữa}: dùng tổng đóng góp của cả hai phân cực:
\begin{equation}
  \text{score}(n,l)_\text{đúng} =
  \underbrace{|V_q^H \mathbf{V}_\text{top}|_{n,l}^2}_{\text{phân cực trên}}
  + \underbrace{|V_q^H \mathbf{V}_\text{bot}|_{n,l}^2}_{\text{phân cực dưới}}
  \label{eq:score_both}
\end{equation}
với $\mathbf{V}_\text{top} = \mathbf{V}_\text{svd}(1:N_1N_2,\,:)$ và
$\mathbf{V}_\text{bot} = \mathbf{V}_\text{svd}(N_1N_2{+}1:\text{end},\,:)$.

%% ============================================================
\section{Chi tiết Cài đặt}
%% ============================================================

\subsection{Tổng quan các thay đổi}

Toàn bộ Mode~B được tích hợp qua \textbf{5 điểm sửa đổi} trong
\texttt{nrPMIReport.m} và \textbf{một hàm mới} (246~dòng).
Không có thay đổi nào trong \texttt{nrCSIReportConfig.m}
(Mode~B tái sử dụng flag \texttt{CodebookMode = 2}
đã có trong cấu trúc Rel-15/16).

\begin{table}[H]
\centering
\caption{Tóm tắt các thay đổi trong \texttt{nrPMIReport.m}}
\begin{tabular}{cll}
\toprule
\# & Vị trí (dòng gần đúng) & Mô tả \\
\midrule
1 & 386\textasciitilde389 & Thêm flag \texttt{isModeB\_r19} \\
2 & 432\textasciitilde444 & Nhánh tính PMI cho Mode~B (bypass toàn bộ luồng cũ) \\
3 & 462 & Bỏ qua xử lý subband cho Mode~B \\
4 & 471\textasciitilde472 & Bỏ qua gán output cho Mode~B (đã gán ở điểm~2) \\
5 & 493\textasciitilde506 & Trong \texttt{getCodebook()}: trả về placeholder thay vì codebook \\
6 & 2401\textasciitilde2517 & Hàm mới \texttt{getPMIType1SinglePanelR19CodebookModeB()} \\
\bottomrule
\end{tabular}
\end{table}

\subsection{Thay đổi 1: Flag \texttt{isModeB\_r19}}

\begin{lstlisting}
% ThangTQ23_128T128R_Rel19: Refined Type I Single-Panel
isType1SinglePanel_r19 = strcmpi(reportConfig.CodebookType,...
                                 'typeI-SinglePanel-r19');
if isType1SinglePanel_r19; isType1SinglePanel = true; end
% modeB all ranks: factored search path (no precomputed codebook)
isModeB_r19 = isType1SinglePanel_r19 && (reportConfig.CodebookMode == 2);
\end{lstlisting}

\noindent\textbf{Lý do}: Flag \texttt{isModeB\_r19} phân nhánh luồng thực thi
\emph{sớm}, trước khi code cũ cố gắng sử dụng codebook tensor
(vốn không tồn tại cho Mode~B).

\subsection{Thay đổi 2: Nhánh tính PMI cho Mode B}

\begin{lstlisting}
if isModeB_r19
    % ThangTQ23_128T128R_Rel19 modeB: factored greedy search, all ranks 1-4
    nVar_scalar = mean(nVar(:), 'omitnan');
    [W, pmiModeB, sinrVec] = getPMIType1SinglePanelR19CodebookModeB( ...
        Hcsirs, reportConfig, nLayers, nVar_scalar, numCSIRSPorts);
    % PMISet: 1-based, i1 = [q1+1, q2+1, n_idx(1)+1, ..., n_idx(v)+1]
    PMISet.i1    = [pmiModeB.q1+1, pmiModeB.q2+1, pmiModeB.n_idx+1];
    PMISet.i2    = pmiModeB.c_idx + 1;
    SINRPerRE    = [];
    SINRPerREPMI = sinrVec(:).';
    SubbandSINRs = sinrVec(:).';
    SINRPerREOut = [];
    codebookOut  = [];
elseif isType2 || isEnhType2
    ...
else
    ...  % Mode A and Rel-15/16
end
\end{lstlisting}

\noindent\textbf{Cấu trúc PMISet}: Mode~B có nhiều hơn 3 chỉ số $i_1$ so với Mode~A
(thêm $n_\text{idx}$ per-layer). Với hạng $v$:
\begin{equation}
  \texttt{i1} = \underbrace{[q_1+1,\; q_2+1]}_{\text{offset}} +
  \underbrace{[n^{(1)}+1,\;\ldots,\;n^{(v)}+1]}_{\text{beam per layer, 1-based}}
  \label{eq:pmiset_i1}
\end{equation}

\subsection{Thay đổi 3 và 4: Bỏ qua các luồng không áp dụng}

\begin{lstlisting}
% Bỏ qua xử lý subband (Mode B chỉ hỗ trợ wideband)
if ~isModeB_r19 && (numSubbands > 1 || (isType2 && ...))
    ...  % subband processing
end

% Form output: Mode B đã gán SINRPerREOut và codebookOut ở trên
if isModeB_r19
    % already set above
elseif isType1SinglePanel || isType1MultiPanel
    SINRPerREOut = SINRPerRE;
    codebookOut  = codebook;
    ...
end
\end{lstlisting}

\subsection{Thay đổi 5: Placeholder trong \texttt{getCodebook()}}

\begin{lstlisting}
if strcmpi(reportConfig.CodebookType,'typeI-SinglePanel-r19')
    if reportConfig.CodebookMode == 1
        % modeA: full precomputed codebook
        codebook = getPMIType1SinglePanelR19CodebookModeA(...);
        [~,~,i2L,i11L,i12L,i13L] = size(codebook);
        indexSetSizes = [i2L,i11L,i12L,i13L];
    else
        % modeB: no precomputed codebook; placeholder passes NaN-check:
        %   ~any(codebook(:)) == false  ->  no early NaN return
        codebook      = complex(1);
        indexSetSizes = [1,1,1,1];
    end
    return;
end
\end{lstlisting}

\noindent\textbf{Lý do kỹ thuật}: Ngay sau \texttt{getCodebook()}, code gốc kiểm tra:
\begin{lstlisting}
if ... (~any(codebook(:)))       % codebook toàn số 0
    PMISet = PMINaNSet;
    return;   % early exit!
end
\end{lstlisting}
Nếu trả về \texttt{codebook = []} hay \texttt{zeros(...)}, hàm sẽ thoát sớm
và trả về NaN. Dùng \texttt{complex(1)} (scalar $1+0j$) vượt qua kiểm tra này
mà không làm sai các luồng code khác.

\subsection{Hàm mới: \texttt{getPMIType1SinglePanelR19CodebookModeB()}}

\subsubsection{Giao diện hàm}

\begin{lstlisting}
function [W, pmiModeB, sinrOut] = getPMIType1SinglePanelR19CodebookModeB(...
    Hcsirs, reportConfig, nLayers, nVar, Pcsirs)
% Inputs:
%   Hcsirs   [nRx x Pcsirs x nRE] - Channel tại CSI-RS RE locations
%   nLayers  scalar                - Rank (1..4)
%   nVar     scalar                - Noise variance (wideband average)
%   Pcsirs   scalar                - Số port CSI-RS (= 128)
% Outputs:
%   W        [Pcsirs x nLayers]   - Ma trận tiền mã hoá tối ưu
%   pmiModeB struct                - {q1, q2, n_idx[1..v], c_idx[1..v]}
%   sinrOut  [nLayers x 1]        - SINR per-layer (ZF, linear)
\end{lstlisting}

\subsubsection{Bước 1: Trung bình Wideband và SVD}

\begin{lstlisting}
% Wideband channel: average over RE dimension
H_wb = mean(Hcsirs, 3);   % [nRx x Pcsirs]

% SVD: leading nLayers right singular vectors
[~, ~, V] = svd(H_wb, 'econ');
V_svd = V(:, 1:nLayers);          % [Pcsirs x nLayers]
V_top = V_svd(1:n_len, :);        % [N1N2 x nLayers]  - phân cực trên
V_bot = V_svd(n_len+1:end, :);    % [N1N2 x nLayers]  - phân cực dưới
\end{lstlisting}

\noindent Vector kỳ dị phải $\mathbf{v}_l$ (cột $l$ của $\mathbf{V}$) nằm trong
không gian hàng của $\mathbf{H}_\text{wb}$ và biểu diễn hướng phát tối ưu cho
tầng $l$. Thuật toán sẽ tìm beam DFT gần nhất với $\mathbf{v}_l$.

\subsubsection{Bước 2: Xây dựng ma trận DFT cho từng offset $(q_1, q_2)$}

\begin{lstlisting}
for q_idx = 0:q_len-1    % q_len = O1*O2 = 16
    q1 = floor(q_idx / O2);
    q2 = mod(q_idx, O2);

    % V_q: [N1N2 x N1N2], cột n = v_{m1,m2} với offset (q1,q2)
    V_q = zeros(n_len, n_len);
    for n_idx = 0:n_len-1
        n_val = N1*N2 - 1 - n_idx;            % Phuong trinh (2)
        n1    = mod(n_val, N1);
        n2    = (n_val - n1) / N1;
        m1    = O1*n1 + q1;                    % Phuong trinh (4)
        m2    = O2*n2 + q2;
        V_q(:, n_idx+1) = getVlm(N1,N2,O1,O2,m1,m2);
    end
    ...
end
\end{lstlisting}

\noindent Ma trận $\mathbf{V}_q \in \mathbb{C}^{N_1 N_2 \times N_1 N_2}$ là
ma trận DFT 2D đã oversampling cho lưới con thứ $q$. Mỗi cột là một ứng viên beam.

\subsubsection{Bước 3: Chọn beam ứng viên}

\begin{lstlisting}
if nLayers == 1
    % Hang 1: vet can tat ca N1*N2 beam (khong gian = Mode A)
    n_cand_matrix = (0:n_len-1)';  % [64 x 1]
else
    % Hang > 1: greedy top-1 beam/tang tu diem chieu ca 2 phan cuc
    scores = abs(V_q' * V_top).^2 + abs(V_q' * V_bot).^2; % [64 x nLayers]
    n_cand = zeros(1, nLayers);
    for l = 1:nLayers
        [~, n_cand(l)] = max(scores(:, l));
        n_cand(l) = n_cand(l) - 1;    % 0-based
    end
    n_cand_matrix = n_cand;  % [1 x nLayers]
end
\end{lstlisting}

\noindent\textbf{Giải thích điểm số} \eqref{eq:score_both}: với beam $v_n$ và
tầng $l$, điểm số đo lường năng lượng kênh được ghi nhận bởi beam $v_n$
khi kết hợp cả hai phân cực anten gNB. Việc thiếu $\mathbf{V}_\text{bot}$
gây chọn beam sai trong $\approx 85\%$ trường hợp trên kênh Rayleigh.

\subsubsection{Bước 4: Tìm kiếm co-phase vét cạn}

\begin{lstlisting}
% Co-phase: 4^nLayers combinations (max 256 for rank 4)
c_combos = dec2base(0:4^nLayers-1, 4, nLayers) - '0';  % [4^v x v]

for n_row = 1:size(n_cand_matrix,1)
    n_try = n_cand_matrix(n_row);  % beam index

    for ci = 1:size(c_combos,1)
        c_vec = c_combos(ci,:);    % [1 x nLayers]

        % Assemble W theo Phuong trinh (1)
        W_top_c = zeros(n_len, nLayers);
        W_bot_c = zeros(n_len, nLayers);
        for l = 1:nLayers
            nb    = n_try(min(l, numel(n_try)));
            n_val = N1*N2 - 1 - nb;
            n1    = mod(n_val, N1);  n2 = (n_val-n1)/N1;
            m1    = O1*n1 + q1;      m2 = O2*n2 + q2;
            v_l   = getVlm(N1,N2,O1,O2,m1,m2);
            W_top_c(:,l) = v_l;
            W_bot_c(:,l) = phi(c_vec(l)) * v_l;
        end
        W_cand = (1/sqrt(nLayers*Pcsirs)) * [W_top_c; W_bot_c];

        % Dung luong Shannon (tieu chi lua chon PMI)
        H_eff = H_wb * W_cand;
        cap = real(log2(det(eye(size(H_eff,1)) + (1/nVar)*(H_eff*H_eff'))));

        if cap > best_cap
            best_cap     = cap;
            best_W       = W_cand;
            best_pmi.q1    = q1;    best_pmi.q2    = q2;
            best_pmi.n_idx = n_try(:)';
            best_pmi.c_idx = c_vec;
        end
    end
end
\end{lstlisting}

\noindent\textbf{Tiêu chí tối ưu}: dung lượng Shannon
$\log_2\det\!\left(\mathbf{I} + \frac{1}{\sigma^2}\mathbf{H}\mathbf{W}\mathbf{W}^H\mathbf{H}^H\right)$
được dùng thay vì SINR đơn giản, vì nó nắm bắt được hiệu quả không gian
của toàn bộ $v$ tầng cùng lúc.

\subsubsection{Bước 5: Tính SINR per-layer (Zero-Forcing)}

\begin{lstlisting}
H_eff = H_wb * W;     % [nRx x nLayers]
W_zf  = pinv(H_eff);  % Zero-forcing combiner
for l = 1:nLayers
    sig     = abs(W_zf(l,:) * H_eff(:,l))^2;         % tin hieu mong muon
    intf    = sum(abs(W_zf(l,:) * H_eff).^2) - sig;  % nhieu lien tang
    sinrOut(l) = sig / (intf + nVar * norm(W_zf(l,:))^2);
end
\end{lstlisting}

%% ============================================================
\section{Kiểm tra và Xác nhận (\texttt{test\_ModeB\_verification.m})}
%% ============================================================

Script kiểm tra \texttt{test\_ModeB\_verification.m} được viết để xác nhận
tính đúng đắn theo 5 tiêu chí độc lập, mỗi tiêu chí kiểm tra một khía cạnh
khác nhau của tiêu chuẩn.

\subsection{TEST 1: Ánh xạ chỉ số}

Kiểm tra trực tiếp các phương trình \eqref{eq:n_mapping}--\eqref{eq:m1m2}:
\begin{itemize}
  \item $i_{1,2,l} = N_1N_2 - 1$, $q=0$ $\Rightarrow$ $m=(0,0)$ $\Rightarrow$ $v_{0,0} = \mathbf{1}$ (\checkmark)
  \item $i_{1,2,l} = 0$, $q=0$ $\Rightarrow$ $m=(O_1(N_1{-}1), O_2(N_2{-}1)) = (60, 12)$ (\checkmark)
  \item Toàn bộ $m_1 \in [0, N_1 O_1{-}1]$ và $m_2 \in [0, N_2 O_2{-}1]$ (\checkmark)
\end{itemize}

\subsection{TEST 2: Chuẩn hoá công suất $\|W\|_F^2 = 1$}

Kiểm tra phương trình \eqref{eq:norm_check} với 4 giá trị ngẫu nhiên của
$(q_1, q_2)$, beam và co-phase, cho mỗi hạng $v \in \{1,2,3,4\}$.
Sai số cho phép $\epsilon < 10^{-10}$.

\subsection{TEST 3: Cấu trúc co-phase $W_\text{bot} = \Phi\,W_\text{top}$}

Kiểm tra trực tiếp rằng mỗi cột của phần dưới $W_\text{bot}$ bằng đúng
$\varphi(c_l)$ lần cột tương ứng của $W_\text{top}$, cho mọi hạng.

\subsection{TEST 4: Phục hồi beam trong kênh LoS}

\noindent\textbf{Thiết lập}: Kênh LoS với beam mục tiêu $m_1=4, m_2=4$
(tương ứng $n_1=1, n_2=1, n=17, i_{1,2,l}=46$, với $q=0$):
\begin{equation}
  \mathbf{H}_\text{LoS} = \mathbf{h}_\text{rx} \cdot v_{4,4}^H
\end{equation}

\noindent\textbf{Bài kiểm tra}: Phép chiếu $\|H_\text{LoS} \cdot [v_n;\, v_n]/\sqrt{2}\|^2$
đạt cực đại tại đúng $i_{1,2,l} = 46$. (\checkmark)

\noindent\textbf{Lỗi đã sửa trong quá trình debug}: Ban đầu dùng transposition thường
$v^\top$ (không liên hợp) thay vì $v^H$ (liên hợp Hermitian):
\begin{lstlisting}
% Sai (bilinear form, khong dat cuc dai tai v=w):
H_los = h_rx * v_target.';

% Dung (Hermitian inner product):
H_los = h_rx * v_target';   % conjugate transpose
\end{lstlisting}
Với vector DFT phức, $v^\top \cdot w \neq v^H \cdot w$, nên bilinear form
không đạt cực đại tại $v = w$.

\subsection{TEST 5: $\text{cap}(\text{Mode~B}) \geq \text{cap}(\text{Mode~A})$}

\noindent\textbf{Thiết lập}: 20 kênh Rayleigh ngẫu nhiên
$\mathbf{H} \in \mathbb{C}^{4 \times 128}$, SNR = 20~dB, hạng~1.
So sánh dung lượng $\log_2\det(\cdot)$ của PMI từ Mode~A và Mode~B.

\noindent\textbf{Kết quả cuối}:

\begin{table}[H]
\centering
\caption{Kết quả TEST 5 sau các sửa chữa}
\begin{tabular}{lcc}
\toprule
Phiên bản & Mode B > Mode A & Mode B = Mode A \\
\midrule
Ban đầu (chỉ $V_\text{top}$, greedy) & 0\% & 15\% \\
Sau sửa scoring ($V_\text{top} + V_\text{bot}$) & 0\% & 15\% \\
Sau vét cạn hạng 1 (64 beam/offset) & 0\% & \textbf{100\%} \\
\bottomrule
\end{tabular}
\end{table}

\noindent Kết quả 100\% bằng nhau xác nhận bổ đề \eqref{eq:equiv_rank1}:
với hạng~1 và tìm kiếm đầy đủ, Mode~B và Mode~A phủ cùng không gian DFT
và luôn tìm được cùng beam tối ưu.

\subsection{Kết quả cuối: Tất cả 5 test PASS}

\begin{tcolorbox}[colback=green!5,colframe=green!50,title=Kết quả kiểm tra Mode B]
\begin{verbatim}
=== Mode B Verification (N1=16, N2=4, O1=4, O2=4, P=128) ===

--- TEST 1: Index mapping ---
  [PASS] i1,2,l=N1N2-1, q=0 -> m=(0,0) -> vlm=all-ones
  [PASS] i1,2,l=0, q=0 -> m=(60,12)=O1*(N1-1), O2*(N2-1)
  [PASS] All m1 in [0,63], m2 in [0,15]

--- TEST 2: Power normalization ||W||_F^2 = 1 ---
  [PASS] rank=1: ||W||_F^2 = 1
  [PASS] rank=2: ||W||_F^2 = 1
  [PASS] rank=3: ||W||_F^2 = 1
  [PASS] rank=4: ||W||_F^2 = 1

--- TEST 3: W structure: W_bot = diag(phi)*W_top ---
  [PASS] rank=1: co-phase structure correct
  [PASS] rank=2: co-phase structure correct
  [PASS] rank=3: co-phase structure correct
  [PASS] rank=4: co-phase structure correct

--- TEST 4: Known rank-1 LoS channel beam recovery ---
  [PASS] LoS beam recovered: i1,2,l=46 -> m=(4,4)

--- TEST 5: cap(Mode B) >= cap(Mode A) [20 random channels] ---
  Mode B > Mode A: 0/20 trials (0%)
  Mode B = Mode A: 20/20 trials (100%)
  cap_B - cap_A:   mean=0.000  min=0.000  max=0.000
  [PASS] cap(Mode B) >= cap(Mode A) in all trials
\end{verbatim}
\end{tcolorbox}

%% ============================================================
\section{Hành trình Debug và các Vấn đề đã Giải quyết}
%% ============================================================

\subsection{Vấn đề 1: MATLAB không cho phép function trong if/else}

\noindent\textbf{Triệu chứng}: Lỗi cú pháp khi cố định nghĩa hàm helper
\texttt{runRI\_and\_PMI()} bên trong khối \texttt{if ThangTQ23\_128T128R\_Rel19}.

\noindent\textbf{Giải pháp}: MATLAB script không cho phép local function
nằm trong khối điều kiện. Giải pháp là \emph{inline} toàn bộ vòng lặp RI
trực tiếp vào thân script.

\subsection{Vấn đề 2: \texttt{complex(1)} placeholder cho NaN-check}

\noindent\textbf{Triệu chứng}: Mode~B trả về NaN mà không thực sự tính PMI.

\noindent\textbf{Nguyên nhân}: Hàm \texttt{getCodebook()} trả về empty array,
kích hoạt kiểm tra \texttt{~any(codebook(:))} $\Rightarrow$ early NaN return.

\noindent\textbf{Giải pháp}: Trả về \texttt{codebook = complex(1)} --- scalar
phức khác không --- để vượt qua điều kiện này. Hàm Mode~B thực sự
được gọi ở nhánh \texttt{if isModeB\_r19} trước khi đến đoạn code đó.

\subsection{Vấn đề 3: Thống nhất hạng 1 và hạng cao}

\noindent\textbf{Thiết kế ban đầu}: Tách hàm \texttt{ModeB\_rank1} (codebook
tính trước) và hàm \texttt{ModeB\_highRank} (greedy search), với
flag \texttt{isModeBHighRank = isModeB\_r19 \&\& nLayers >= 2}.

\noindent\textbf{Quyết định thay đổi}: Người dùng yêu cầu thiết kế nhất quán:
một hàm duy nhất cho mọi hạng, tên \texttt{getPMIType1SinglePanelR19CodebookModeB}
đồng nhất với \texttt{getPMIType1SinglePanelR19CodebookModeA}.

\noindent\textbf{Kết quả}: Hạng~1 trong hàm thống nhất sử dụng vét cạn nên
đạt cùng kết quả tối ưu như hàm riêng, đồng thời code gọn hơn.

\subsection{Vấn đề 4: Điểm số beam thiếu phân cực dưới}

\noindent\textbf{Triệu chứng}: TEST~5 fail, Mode~B thua Mode~A trong 85\% thử nghiệm.

\noindent\textbf{Phân tích}: Code gốc dùng
$\text{scores} = |V_q^H \mathbf{V}_\text{top}|^2$ (chỉ phân cực trên).
Với kênh Rayleigh, $\mathbf{H}_\text{top}$ và $\mathbf{H}_\text{bot}$
độc lập, nên SVD vector $\mathbf{v}_l \in \mathbb{C}^{128}$ phân bổ
năng lượng đồng đều giữa hai nửa. Bỏ qua nửa dưới làm mất 50\% thông tin
về hướng beam tối ưu.

\noindent\textbf{Sửa chữa}: Cộng thêm đóng góp từ $\mathbf{V}_\text{bot}$
(phương trình \eqref{eq:score_both}). Sau đó TEST~5 vẫn fail vì vấn đề
hạng~1 (vét cạn vs. greedy top-1) là nguyên nhân chính.

\subsection{Vấn đề 5: Hạng 1 cần vét cạn, không phải greedy top-1}

\noindent\textbf{Phân tích}: Với $q$ cố định, greedy chỉ chọn 1 beam tốt nhất
theo điểm SVD, rồi thử 4 co-phase. Tổng: $16 \times 4 = 64$ đánh giá.
Mode~A đánh giá $16 \times 64 \times 4 = 4{,}096$ -- nhiều hơn 64 lần.
Beam tốt nhất theo điểm SVD \emph{không nhất thiết} là beam cho dung lượng
cao nhất sau khi tối ưu co-phase.

\noindent\textbf{Sửa chữa}: Khi \texttt{nLayers == 1},
\texttt{n\_cand\_matrix = (0:n\_len-1)'} (tất cả 64 beam) thay vì chỉ top-1.
Điều này khôi phục vét cạn đầy đủ và đảm bảo kết quả tối ưu cho hạng~1,
đồng thời không làm chậm đáng kể ở hạng cao hơn.

%% ============================================================
\section{Bảng kiểm tra Tuân thủ Tiêu chuẩn (Mode B)}
%% ============================================================

\begin{table}[H]
\centering
\caption{Mode B: Kiểm tra tuân thủ TS~38.214 \S5.2.2.2.1a}
\begin{tabular}{p{7cm}p{5.5cm}c}
\toprule
Yêu cầu tiêu chuẩn & Cài đặt & Trạng thái \\
\midrule
Bảng 5.2.2.2.1a-6/7: ánh xạ $n = N_1N_2-1-i_{1,2,l}$
  & \texttt{n\_val = N1*N2-1-n\_idx} & \checkmark \\
$n_1 = n \bmod N_1$, $n_2 = \lfloor n/N_1 \rfloor$
  & \texttt{mod(n\_val,N1)}, \texttt{(n\_val-n1)/N1} & \checkmark \\
$m_1 = O_1 n_1 + q_1$, $m_2 = O_2 n_2 + q_2$
  & \texttt{O1*n1+q1}, \texttt{O2*n2+q2} & \checkmark \\
$q_1 \in \{0,\ldots,O_1{-}1\}$, $q_2 \in \{0,\ldots,O_2{-}1\}$
  & \texttt{floor(q\_idx/O2)}, \texttt{mod(q\_idx,O2)} & \checkmark \\
$i_{1,2,l} \in \{0,\ldots,N_1N_2{-}1\}$ per layer
  & \texttt{n\_cand in 0:N1N2-1} & \checkmark \\
Cấu trúc $W = [V; \Phi V]/\sqrt{vP}$
  & \texttt{[W\_top\_c; W\_bot\_c]/sqrt(v*P)} & \checkmark \\
$\varphi(c) = e^{j\pi c/2}$, $c \in \{0,1,2,3\}$
  & \texttt{phi = @(x) exp(1i*pi*x/2)} & \checkmark \\
Chuẩn hoá $\|W\|_F^2 = 1$
  & Đã kiểm tra TEST 2 & \checkmark \\
Hạng 1: beam chọn tối ưu (Mode B $\geq$ Mode A)
  & Đã kiểm tra TEST 5 (20/20 trials) & \checkmark \\
\midrule
\multicolumn{2}{l}{\textit{Hạn chế đã biết}} & \\
Subband PMI cho Mode B
  & Chưa cài đặt (wideband only) & -- \\
Hạng 2--4: tối ưu toàn cục
  & Greedy (xấp xỉ) & \textasciitilde \\
\bottomrule
\end{tabular}
\end{table}

%% ============================================================
\section{Kết luận}
%% ============================================================

Mode~B của Refined Codebook Type~I Single-Panel (TS~38.214 \S5.2.2.2.1a)
đã được cài đặt thành công vào thư viện MATLAB tùy chỉnh với các
đặc điểm sau:

\begin{enumerate}
  \item \textbf{Đúng tiêu chuẩn}: Ánh xạ chỉ số, cấu trúc bộ tiền mã hoá,
        và chuẩn hoá công suất đều tuân thủ chính xác Bảng~5.2.2.2.1a-7.

  \item \textbf{Tối ưu hạng~1}: Sử dụng vét cạn đầy đủ ($16 \times 64 \times 4 = 4{,}096$
        ứng viên), đảm bảo Mode~B $\geq$ Mode~A (kiểm chứng 20/20 trials).

  \item \textbf{Khả thi bộ nhớ hạng cao}: Greedy Factored Search giữ số
        đánh giá trong $O(O_1 O_2 \times N_1 N_2 \times 4^v)$ thay vì
        $O((N_1 N_2)^v \times 4^v)$, tránh yêu cầu bộ nhớ GB.

  \item \textbf{Tích hợp sạch}: 5 điểm sửa đổi nhỏ trong \texttt{nrPMIReport.m},
        không ảnh hưởng đến Mode~A và các luồng Rel-15/16 hiện có.
\end{enumerate}

\noindent\textbf{Công việc tiếp theo đề xuất:}
\begin{itemize}
  \item Chạy so sánh hiệu năng Mode~A vs Mode~B trên kênh CDL-B/C
        cho hạng 2, 3, 4 để đánh giá mức độ cải thiện thực tế.
  \item Cải thiện greedy hạng cao bằng deflation (loại bỏ đóng góp beam đã chọn
        trước khi chọn beam tiếp theo).
  \item Mở rộng sang subband PMI cho Mode~B.
\end{itemize}

\end{document}
